% ===================================================================
%  TÁBLA Uimh. 14 — An tSráid. — Na Ceannaithe.
%  Traduction 2026 d'apres texto/io, controlee sur texto/fr, texto/en
%  et texto/es. Consigne : voir texto/ga/10-tabla-01.tex.
%
%  LE TABLEAU DES METIERS, ET C'EST LE PLUS LONG DE LA COLONNE. Le
%  partage y est plus fin qu'ailleurs, parce que les metiers ne sont
%  pas tous du meme age. Ce qui se fait de ses mains depuis toujours
%  garde son mot irlandais — « gabha » le forgeron, « táilliúir » le
%  tailleur, « bearbóir » le barbier, « mangaire » le colporteur,
%  « ceannaí » le marchand, « fuarán » la fontaine, « cónra » le
%  cercueil. Ce qui s'est invente au XIXe siecle arrive avec son nom :
%  « grianghrafadóir », « clóscríobhán », « liteagrafaí »,
%  « rabhadán », « omnaibus », « fiacra », « consal ». Le lieu de
%  l'apprentissage decide ici comme ailleurs, mais il se double d'une
%  DATE : le mot suit l'objet, et l'objet a un age.
%
%  ZERO INVERSION. Le tableau enumere autant que le tableau 12, et
%  pour la meme raison il ne coute rien. Les deux passages qui
%  auraient pu couter se rangent seuls : « la skribo-chambro (32) di
%  mea kuzulo (34) » — genitif postpose, comme l'ido — et « avan la
%  regimento (41) marchis la tamburestro (42) », ou l'irlandais
%  antepose son complement de lieu et laisse le VSO derriere, ce qui
%  rend exactement l'ordre de l'ido.
% ===================================================================

%%K t14-apar-1 apar
\VUsaut{4.0mm}
\VUtitre{15.0pt}{60}{TÁBLA Uimh. 14}
\VUsaut{3.0mm}

%%K t14-tit-1 sub
\VUpk{11.4pt}{40}{An tSráid. --- Na Ceannaithe.}
\VUsaut{3.0mm}

%%K t14-01-1 p
1. --- Tháinig mo chara Seán, a bhfuil cónaí air faoin tuath, chun trí lá a chaitheamh
i dteannta mo mhuintire. \VUgras{Go dtí sin} ní raibh caoi aige cathair mhór a
fheiceáil; dá bhrí sin caithimid ár laethanta go léir ag siúl, agus mínímse dó an rud
nach eol dó.

%%K t14-02-1 p
2. --- Ar dtús chuamar ag féachaint ar an \VUgras{réadlann} \textsuperscript{(1)}, a
chuir an-suim ann. Ag filleadh dúinn tríd an \VUgras{tsráid} \textsuperscript{(3)} ina
bhfuil na \VUgras{Folcadáin Thurcacha} \textsuperscript{(2)}, casadh \VUgras{sochraid}
\textsuperscript{(4)} orainn. Os comhair \VUgras{mhórshiúl na sochraide} shiúil
\VUgras{an chléir,} rompu \VUgras{cóiste na marbh,} ina raibh an \VUgras{chónra}
curtha. Bhí a lán cairde i dteannta an teaghlaigh a bhí \VUgras{faoi bhrón}.

%%K t14-02-2 p
Nuair a bhí an mhuintir a lean tar éis dul thart, thrasnaíomar an tsráid os comhair an
\VUgras{tí bheorach} mhóir \textsuperscript{(5)} ina raibh a lán \VUgras{custaiméirí}
ag ól \VUgras{beorach} ar an teirís, faoi scáth an \VUgras{scáthláin ghloine}
\textsuperscript{(6)}.

%%K t14-03-1 p
3. --- Bhí ionadh an domhain ar Sheán faoin \VUgras{sciath} a bhí curtha idir siopa an
\VUgras{díoltóra biotáille} \textsuperscript{(7)} agus siopa an \VUgras{díoltóra
ceoil} \textsuperscript{(10)}. D'fhiafraigh sé díom cad ba bhrí leis; d'fhreagair mé
gur \VUgras{consalacht} \textsuperscript{(8)} Shasana a chomharthaigh sé, agus
tharraing mé a aird ar an \VUgras{gconsal} \textsuperscript{(9)} a bhí ar an mbalcóin.

%%K t14-tit-2 sub
\VUpk{10.2pt}{40}{Turas na siopaí.}
\VUsaut{3.0mm}

%%K t14-04-1 p
4. --- Ar ndóigh stadamar os comhair thaispeántas na \VUgras{n-uirlisí ceoil,}
\VUgras{na n-harmóiniamaí,} \VUgras{na n-orgán} agus \VUgras{na bpianónna}; d'fhéach
muid ar chlúdaigh mhaisithe na \VUgras{scórtha} agus na \VUgras{bpíosaí ceoil} don
phianó, agus bhí meas againn ar an \VUgras{lir} adhmaid órtha a úsáidtear mar
chomhartha.

%%K t14-05-1 p
5. --- Chuir na scamaill deannaigh a rinne na \VUgras{scuabadóirí} \textsuperscript{(11)}
agus an \VUgras{carr scuabtha} \textsuperscript{(12)} orainn dul go tapa thar
\VUgras{throscán} breá an \VUgras{díoltóra troscáin} \textsuperscript{(13)} agus thar
thaispeántas \VUgras{sorn} agus \VUgras{lampaí} an \VUgras{stánadóra}
\textsuperscript{(14)}.

%%K t14-06-1 p
6. --- Chomh maith leis sin b'éigean dúinn ag an áit sin an cosán a fhágáil, mar bhí
sé lán de phacáistí a tógadh as \VUgras{veain iompair} \textsuperscript{(16)} chun iad
a chur i \VUgras{siopa} \textsuperscript{(15)} a bhí díreach tógtha ar cíos.

%%K t14-06-2 p
Chuir sin cosc orainn carranna breátha an \VUgras{déantóra carranna}
\textsuperscript{(17)} a fheiceáil, ach níor chuir sé cosc orainn stad chun féachaint
ar an \VUgras{bpacálaí} \textsuperscript{(19)} agus é ag déanamh ciste dá chomharsa,
an \VUgras{díoltóir rothar agus gluaisteán} \textsuperscript{(18)}. Ansin, ó bhí sé
beagán déanach cheana, d'fhilleamar abhaile.

%%K t14-tit-3 sub
\VUpk{10.2pt}{40}{Siúlóid tríd an mbaile.}
\VUsaut{3.0mm}

%%K t14-07-1 p
7. --- Lá arna mhárach mhol mo mháthair do Sheán go dtógfaí a phictiúr, agus chuir sí
de chúram orm dul in éineacht leis chuig an \VUgras{ngrianghrafadóir}
\textsuperscript{(20)}. Tar éis an lóin chuamar isteach i \VUgras{gceardlann} an
ealaíontóra, mar ar chaith muid tamall maith ag feitheamh. Chun sinn féin a
choimeád gnóthach, d'fhéach muid ar na \VUgras{portráidí} \textsuperscript{(22)} a bhí
ar crochadh ar an mballa.

%%K t14-07-2 p
Nuair a tháinig ár seal, níor mhair an obair i bhfad, agus chomh luath is a chríochnaigh
mo chara a sheasamh os comhair an \VUgras{cheamara} \textsuperscript{(21)}, chuamar
síos.

%%K t14-08-1 p
8. --- Ar an gcéad urlár chonaiceamar trí dhoras an \VUgras{bhearbóra}
\textsuperscript{(23)}, a bhí ar oscailt, a bhuachaill ag bearradh gruaige custaiméara.

%%K t14-08-2 p
Nuair a shroicheamar an tsráid ghabhamar thar an \VUgras{siopa tobac}
\textsuperscript{(24)}. Bhí Seán chomh tógtha sin leis an \VUgras{laindéar} dearg
\textsuperscript{(25)} agus leis an \VUgras{bpíopa mór} a úsáidtear mar
\VUgras{chomhartha} \textsuperscript{(26)}, gur bhuail sé faoi bheirt sheanuaisle a bhí
ag comhrá agus \VUgras{snaoisín á chaitheamh acu} \textsuperscript{(27)}.

%%K t14-tit-4 sub
\VUpk{10.2pt}{40}{An Siopa Cácaí.}
\VUsaut{3.0mm}

%%K t14-09-1 p
9. --- Tharraing na \VUgras{nótaí bainc} agus na \VUgras{boinn airgid} as gach tír a
bhí ar taispeáint i \VUgras{bhfuinneog} an \VUgras{fhir mhalartaithe airgid}
\textsuperscript{(29)} ár n-aird tamall, ach ó bhí sé in am greim a ithe, chuamar
isteach chuig an \VUgras{mbáicéir cácaí} \textsuperscript{(31)} gan stad níos faide.

%%K t14-10-1 p
10. --- Nuair a chonaiceamar na \VUgras{sólaistí} go léir a bhí sa siopa,
\VUgras{cácaí, cácaí meala, brioscaí} agus gach cineál \VUgras{milseán,} níorbh
fhurasta dúinn rogha a dhéanamh. Faoi dheireadh roghnaigh Seán na \VUgras{cácaí
uachtair,} agus níor thóg mise ach \VUgras{toirtín} beag \VUgras{silíní,} mar
\VUgras{milleann an tsiúcra mo chuid fiacla,} agus ní theastaíonn uaim ar chor ar bith
a bheith ag dul rómhinic chuig an \VUgras{bhfiaclóir} \textsuperscript{(30)}.

%%K t14-tit-5 sub
\VUpk{10.2pt}{40}{An Clóscríobh.}
\VUsaut{3.0mm}

%%K t14-11-1 p
11. --- Chun \VUgras{clóscríobhán} a thaispeáint do mo chara, ceann ar chuala sé trácht
air ach nach raibh aithne aige air, chuamar isteach in \VUgras{oifig}
\textsuperscript{(32)} mo \VUgras{chol ceathrair} \textsuperscript{(34)}. Fuaireamar é
ag deachtú a chuid litreacha dá \VUgras{rúnaí} \textsuperscript{(35)}, ar
\VUgras{gearrscríobhaí} é. Ar éigean a bhí litir críochnaithe nuair a rinne
\VUgras{clóscríobhaí} \textsuperscript{(33)} cóip di ar a inneall. Ó nár theastaigh
uainn cur isteach ar chúraimí mo ghaoil, níor fhan muid aige ach cúpla nóiméad.

%%K t14-12-1 p
12. --- Faoi oifig mo chol ceathrair tá \VUgras{uaireadóirí-seodóir} mór
\textsuperscript{(37)} a bhfuil ceird an ghaibhne óir aige chomh maith. Tá a lán
\VUgras{clog, clog balla, uaireadóirí, slabhraíní} agus \VUgras{seoda} luachmhara
ina shiopa iontach, mar shampla \VUgras{muinceáin phéarlaí,} \VUgras{bráisléid} óir,
\VUgras{fáinní cluaise} agus \VUgras{fáinní méire} maisithe le \VUgras{clocha} agus
le \VUgras{diamaint}. Tá ceann de na fuinneoga curtha ar leataobh do na \VUgras{hearraí
óir} agus tá bailiúchán mór \VUgras{sceanra} airgid ann.

%%K t14-13-1 p
13. --- Ag casadh dúinn ag coirnéal na sráide thug mé faoi deara go raibh an
\VUgras{pláta} \textsuperscript{(36)} in aice le \VUgras{huimhir} an tí athraithe. Bhí
mé ar tí mo bharúil a rá os ard nuair a chualathas fuaimeanna aoibhne ó \VUgras{cheol}
an airm. Thosaíomar ag rith i gcoinne na reisiminte, gan cuimhneamh go raibh
\VUgras{sí} le gabháil thar shiopaí an \VUgras{díoltóra hataí} \textsuperscript{(39)}
agus an \VUgras{díoltóra lámhainní} \textsuperscript{(40)}, agus go mbeimis chomh
maith céanna as chun a mórshiúl a fheiceáil dá bhfanaimis os comhair fhuinneog an
\VUgras{radharceolaí} \textsuperscript{(38)}, atá lán de \VUgras{phionsnéis,}
\VUgras{spéaclaí} agus \VUgras{ghloiní beaga}.

%%K t14-tit-6 sub
\VUpk{10.2pt}{40}{An Mórshiúl.}
\VUsaut{3.0mm}

%%K t14-14-1 p
14. --- Os comhair na \VUgras{reisiminte} \textsuperscript{(41)} shiúil \VUgras{máistir
na ndrumaí} \textsuperscript{(42)}, agus ina dhiaidh na \VUgras{drumadóirí}
\textsuperscript{(43)}, na \VUgras{stocairí} \textsuperscript{(44)} agus an
\VUgras{cheolfhoireann} ar fad. Bhí \VUgras{an ginearál} \textsuperscript{(45)}, a bhí
ar an gcearnóg lena aidiúnach, stoptha in aice le \VUgras{scaird uisce}
\textsuperscript{(49)} an \VUgras{fhuaráin} \textsuperscript{(48)} chun féachaint ar an
\VUgras{mórshiúl}.

%%K t14-14-2 p
Bheannaigh \VUgras{oifigigh na foirne} \textsuperscript{(46)}, an \VUgras{coirnéal}
agus an \VUgras{leas-choirnéal} dó lena gclaíomh. Bhí \VUgras{na maoir} os comhair a
gcuid \VUgras{cathlán} \textsuperscript{(47)} agus bhí gach \VUgras{captaen} i gceannas
a \VUgras{chomplachta,} fad is a bhí na \VUgras{leifteanaint,} na
\VUgras{fo-leifteanaint} agus na \VUgras{fo-oifigigh} ina n-áit ghnách in aice lena
gcuid fear.

%%K t14-15-1 p
15. --- Bhí a lán daoine tar éis cruinniú ar an \VUgras{gcrosbhóthar}
\textsuperscript{(50)}; tháinig na ceannaithe, an \VUgras{díoltóir scáthanna báistí}
\textsuperscript{(51)}, \VUgras{fear an stánaithe} \textsuperscript{(53)} agus an
\VUgras{gunnadóir} \textsuperscript{(54)} amach chun na \VUgras{saighdiúirí} a
fheiceáil. Chuir na feithiclí go léir, \VUgras{fiacraí} \textsuperscript{(52)},
\VUgras{carranna crainn} \textsuperscript{(57)} agus fiú an \VUgras{bairille
uiscithe} \textsuperscript{(56)}, iad féin ar leataobh le hais an chosáin.

%%K t14-tit-7 sub
\VUpk{10.2pt}{40}{Radharcanna ar an tSráid.}
\VUsaut{3.0mm}

%%K t14-15-2 p
Thrasnaigh \VUgras{taistealaí tráchtála} \textsuperscript{(55)} an tsráid go mear agus
a \VUgras{bhoscaí samplaí} ar iompar ina láimh aige, agus d'iompaigh na daoine a bhí
ina suí ar \VUgras{bharr} \textsuperscript{(59)} an \VUgras{omnaibus}
\textsuperscript{(58)} thart chun sult a bhaint as an radharc. B'éigean d'\VUgras{amhránaí
fáin} bocht \textsuperscript{(60)} agus a \VUgras{orgán sráide} \textsuperscript{(61)}
aige a \VUgras{amhrán} a bhriseadh, mar bhí a ghuth báite ag torann na ndrumaí.

%%K t14-16-1 p
16. --- Nuair a shroich an reisimint an \VUgras{siopa éadaigh}
\textsuperscript{(64)}, d'fhág na \VUgras{mná fuála} \textsuperscript{(63)} agus na
\VUgras{táilliúirí} \textsuperscript{(62)} a gcuid \VUgras{inneall fuála} agus a
gcuid oibre chun féachaint amach an fhuinneog. B'éigean don \VUgras{cheannaí}
\textsuperscript{(65)}, a bhí díreach tar éis \VUgras{cultacha} nua
\textsuperscript{(66)} a chur ar a \VUgras{thaispeántas,} na \VUgras{brait}
\textsuperscript{(67)} agus na \VUgras{línéadaí} a bhí curtha amach ar an gcosán, in
aice le \VUgras{béal an draenacháin} \textsuperscript{(68)}, a thabhairt isteach ar
eagla go leagfadh an slua iad; b'éigean fiú do \VUgras{mhangaire}
\textsuperscript{(69)} aosta, a raibh bean an dorais ag margáil leis faoi stocaí, dul
isteach i bpasáiste chun a dhíolachán a chríochnú.

%%K t14-16-2 p
Chuamar in éineacht leis an reisimint chomh fada leis an mbeairic, \VUgras{agus,} agus
sinn an-sásta lenár lá, d'fhilleamar in am an dinnéir.

%%K t14-tit-8 sub
\VUpk{10.2pt}{40}{Ceirdeanna agus Gairmeacha Éagsúla.}
\VUsaut{3.0mm}

%%K t14-17-1 p
17. --- An lá dár gcionn mhol m'athair dúinn cuairt a thabhairt ar chlólann nuachtáin
na háite. Ghlacamar leis go fonnmhar.

%%K t14-17-2 p
Is \VUgras{clódóir} é an fear céanna agus is \VUgras{leabharcheannaí}
\textsuperscript{(70)} \VUgras{(= díoltóir leabhar),} \VUgras{foilsitheoir,}
\VUgras{díoltóir páipéir} agus \VUgras{ceanglóir} é chomh maith. Chuir sé ar an gcéad
urlár \VUgras{seomra eagarthóireachta} \textsuperscript{(71)} an nuachtáin, agus an
\VUgras{cheardlann ghreanta} freisin, ina n-oibríonn na \VUgras{liteagrafaithe}
\textsuperscript{(72)} agus na \VUgras{greanadóirí copair} \textsuperscript{(73)}, agus
ar deireadh an \VUgras{seomra cló,} ina bhfuil na \VUgras{hoibrithe cló}
\textsuperscript{(74)}. Tá an \VUgras{chlólann} féin \textsuperscript{(75)}, na
\VUgras{preasanna cló} agus na hinnill éagsúla ar urlár na talún.

%%K t14-tit-9 sub
\VUpk{10.2pt}{40}{Cuireann Seán a lán ceisteanna.}
\VUsaut{3.0mm}

%%K t14-18-1 p
18. --- Ag fágáil na clólainne dúinn stad Seán agus ionadh mór air os comhair bosca
beag gloine a bhí curtha ar cholún, le hais an \VUgras{tsiopa mionearraí}
\textsuperscript{(77)}, agus d'fhiafraigh sé cad chuige a n-úsáidtear é. Mhínigh
m'athair dó gur \VUgras{rabhadán dóiteáin} \textsuperscript{(76)} atá ann. Chomh maith
leis sin bhí gach rud sa chathair ina ábhar iontais do mo chara, ón
\VUgras{bhfuarán cloiche} simplí \textsuperscript{(78)} agus ón \VUgras{tríchruthach
iompair} éadrom \textsuperscript{(80)} a bhí á thiomáint ag \VUgras{gasúr}
\textsuperscript{(79)} faoi éide, go dtí \VUgras{carr an phoist}
\textsuperscript{(81)}.

%%K t14-19-1 p
19. --- Fad is a d'fhág m'athair sinn tamall chun labhairt le \VUgras{bailitheoir na
gcánacha} \textsuperscript{(83)}, a bhí stoptha ag comhrá le \VUgras{habhcóide}
\textsuperscript{(82)} agus le \VUgras{nótaire} \textsuperscript{(84)}, thugamar
caitheamh aimsire dúinn féin ag leanúint le súil an tsiúil ar an tsráid. Ghabh na
daoine ba éagsúla tharainn: tháinig \VUgras{buachaill an bháicéara}
\textsuperscript{(85)}, agus \VUgras{pastae} iontach á iompar aige i gciseán, taobh
thiar de \VUgras{dhíoltóir éadaí} \textsuperscript{(86)}; bhí \VUgras{lasadóir na
lampaí} \textsuperscript{(87)} ag comhrá le \VUgras{fear an gháis}
\textsuperscript{(88)}; bhí \VUgras{bean rialta} \textsuperscript{(89)} agus
dílleachtaí ar láimh aici ag filleadh ar a clochar.

%%K t14-tit-10 sub
\VUpk{10.2pt}{40}{Ar ár mbealach abhaile.}
\VUsaut{3.0mm}

%%K t14-20-1 p
20. --- An nóiméad ar tháinig m'athair suas linn, rinne Seán gáire mór nuair a chonaic
sé aghaidh shalach agus éadaí aisteacha beirte \VUgras{ghlantóirí simléar}
\textsuperscript{(91)} a bhí clúdaithe le súiche.

%%K t14-21-1 p
21. --- Theastaigh uaim planda a cheannach ó \VUgras{bhean na mbláth}
\textsuperscript{(92)} do mo mháthair, ach thug m'athair le fios dom go mbeadh sé
róthrom le hiompar agus gurbh fhearr \VUgras{oráistí} a fháil. Chuamar i ngar don
\VUgras{bhean díolta} \textsuperscript{(94)}, agus nuair a chríochnaigh an
\VUgras{bhanoide} \textsuperscript{(93)}, a bhí ag roghnú dá dalta, a ceannach,
thugamar orainn féin freastal a fháil.

%%K t14-22-1 p
22. --- Ag filleadh abhaile dúinn casadh roinnt aithne orainn: seanchara de chuid mo
mhuintire, a bhí ag ligean a huillinne ar a \VUgras{bean chompánach}
\textsuperscript{(95)}, \VUgras{an t-innealtóir} \textsuperscript{(96)} droichead
agus bóithre, agus duine de mo chol ceathracha lena \VUgras{buime}
\textsuperscript{(98)}.

%%K t14-22-2 p
Ansin casadh \VUgras{bean hataí} \textsuperscript{(97)} mo mháthar orainn, agus beirt
\VUgras{mhanach} \textsuperscript{(99)} nach raibh eolas acu ar an gcathair, a tháinig
chugainn chun ceist a chur ar m'athair faoin mbealach chun an stáisiúin.
